% Part: sets-functions-relations
% Chapter: arithmetization
% Section: reflections

\documentclass[../../../include/open-logic-section]{subfiles}

\begin{document}

\olfileid{sfr}{arith}{ref}
\olsection{కొన్ని తాత్త్విక పరామర్శలు}

ఇంతవరకు సాంకేతిక వివరాలు చూశాం. కానీ అవి ఏమి సాధించాయి?

వాటి ద్వారా మనకు కొంత {అందమైన} శుద్ధ గణితం లభించిందనేది దాదాపు
వివాదరహితం. అంతేకాక కొన్ని లోతైన భావనాత్మక విజయాలు దక్కాయి. వాస్తవ
సంఖ్యలకు, కరణీయ సంఖ్యలకు మధ్య కీలకమైన తేడాను సంపూర్ణతా ధర్మం
వ్యక్తపరుస్తుందని గ్రహించడం గొప్ప అంతర్దృష్టి. ఇంకా, వాస్తవ సంఖ్యలను
డెడెకిండ్ కోతలుగా స్పష్టంగా నిర్మించడం విశ్లేషణ విషయానికి దృఢమైన
పునాది ఇస్తుంది. \emph{సంపూర్ణ క్రమిత క్షేత్రం} అనే భావన సుసంగతమని
మనకు తెలుసు; ఎందుకంటే కోతలు సరిగ్గా అటువంటి క్షేత్రాన్నే ఏర్పరుస్తాయి.

అయినప్పటికీ, ఈ విజయాల గురించి కొన్ని సందేహాలను వెల్లడించాలి.

మొదట, వాస్తవ సంఖ్యలను కోతలుగా భావించడం, వాటి పరిచితమైన (బహుశా
అనంతమైన) దశాంశ విస్తరణలుగా భావించడం కంటే \emph{మరింత} కచ్చితమని
స్పష్టం కాదు. దశాంశాల ద్వారా చేసే ఈ రెండవ ``నిర్మాణం'', కౌషీ క్రమాల
ద్వారా వాస్తవ సంఖ్యల నిర్మాణాన్ని కొంత పోలి ఉంటుంది; నిజానికి అది
పదిహేడవ శతాబ్దపు ఆరంభం నుంచే గణితవేత్తలకు సారాంశంగా తెలుసు
(\olref[cauchy]{sec} చూడండి). వాస్తవ సంఖ్యలకు సంపూర్ణతా ధర్మం ఉందని
గ్రహించడమే కచ్చితత్వంలో నిజమైన పెరుగుదల; వాస్తవ సంఖ్యలను కొన్ని
ప్రత్యేక సమితులుగా నిర్మించగలగడం ఒక్కటే అంత ఆసక్తికరమైన విషయం కాకపోవచ్చు.

(ఇంకా సులభమైన) పూర్ణ సంఖ్యల లేదా కరణీయ సంఖ్యల అంకగణితీకరణ ఆయా
విషయాల్లో కచ్చితత్వాన్ని పెంచుతుందా అనేది మరింత అస్పష్టం. ఇక్కడ ఒక
సులభమైన పరిశీలన ఉపయుక్తం. సహజ సంఖ్యల క్రమయుగ్మాల తుల్యతా వర్గాలుగా
పూర్ణ సంఖ్యలను \emph{నిర్మించి}, తరువాత పూర్ణ సంఖ్యల క్రమయుగ్మాల
తుల్యతా వర్గాలుగా కరణీయ సంఖ్యలను నిర్మించి, తరువాత కరణీయ సంఖ్యల
సమితులుగా వాస్తవ సంఖ్యలను నిర్మించిన వెంటనే, ఆ నిర్మాణాలను మనం
\emph{మర్చిపోతాం}. ముఖ్యంగా, గణిత నిరూపణలో ఈ నిర్మాణాలను ఎవరూ
\emph{ప్రయోగించాలనుకోరు} (ఆ నిర్మాణాలు అనుకున్నట్లు ప్రవర్తిస్తాయని
నిరూపించే సందర్భం తప్ప). సహజ సంఖ్యల సమితుల సమితుల సమితుల సమితుల
సమితుల సమితుల సమితి గురించి మాట్లాడటంకంటే, వాస్తవ సంఖ్య గురించి
నేరుగా మాట్లాడటం చాలా సులభం.
%(ఎందుకంటే వాస్తవ సంఖ్యలను కరణీయ సంఖ్యల సమితులుగా నిర్మించాం; కరణీయ సంఖ్యలను పూర్ణ సంఖ్యల క్రమయుగ్మాల తుల్యతా వర్గాలుగా, అంటే పూర్ణ సంఖ్యల సమితుల సమితుల సమితులుగా నిర్మించాం; మొదలైనవి.)
%2/3 = [2,3]
%	అంటే {<2,3>, ... }
%	{ {{2},{2,3}}, ... }
%
%వాస్తవ సంఖ్యలు
%	కరణీయ సంఖ్యల సమితులు
%
%కరణీయ సంఖ్యలు
%	పూర్ణ సంఖ్యల సమితుల సమితుల సమితులు
%
%పూర్ణ సంఖ్యలు
%	సహజ సంఖ్యల సమితుల సమితుల సమితులు

ఈ నిర్వచనాలు, అధిభౌతికంగా చెప్పాలంటే, పూర్ణ సంఖ్యలు, కరణీయ సంఖ్యలు,
వాస్తవ సంఖ్యలు \emph{ఏమిటో} మనకు చెబుతాయనేది అన్నిటికన్నా ఎక్కువ
సందేహాస్పదం. అంటే, వాస్తవ సంఖ్యలు (ఉదాహరణకు) కొన్ని (సమితుల సమితుల
సమితుల\ldots) సమితులే అని అనుకోవడం సందేహాస్పదం. అటువంటి అభిప్రాయానికి
ప్రధాన అడ్డంకి ఏమిటంటే, ఈ నిర్మాణాన్ని అనేక భిన్నమైన మార్గాల్లో
చేయవచ్చు. వాస్తవ సంఖ్యల విషయంలో నిజంగానే ఆసక్తికరంగా భిన్నమైన కొన్ని
నిర్మాణాలు ఉన్నాయి (\olref[cauchy]{sec} చూడండి). అయితే భిన్న నిర్మాణాలను
పొందే అతి సాధారణమైన మార్గం ఇది: \olref[sfr][rel][ref]{sec}లో
చేసినట్లుగా, క్రమయుగ్మాలను కొద్దిగా భిన్నంగా నిర్వచించవచ్చు; ఆ ప్రత్యామ్నాయ
క్రమయుగ్మ భావనను వాడి ఉంటే, మన నిర్మాణాలు ఇంతే చక్కగా పనిచేసేవి,
కానీ చివరికి భిన్నమైన వస్తువులు వచ్చేవి. అందువల్ల పూర్ణ, కరణీయ,
వాస్తవ సంఖ్యలకు అనేక పోటీ సమితి-సిద్ధాంత నిర్మాణాలు ఉన్నాయి. ఇప్పుడు
పూర్ణ సంఖ్యలు (ఉదాహరణకు) \emph{ఆ} సమితులు కాకుండా \emph{ఈ} సమితులే
అని చెప్పడం యాదృచ్ఛికంగా (మరియు ఇబ్బందికరంగా) ఉంటుంది.
(\olref[sfr][rel][ref]{sec}లో మాదిరిగా, ఇది
\citealt{Benacerraf1965} ప్రసిద్ధి చేసిన వాదనకు ఒక ఉదాహరణ.)

ఇంకొక అంశాన్ని ప్రస్తావించాలి: మన నిర్మాణాల్లో కొంత \emph{వింత} ఉంది.
మనం సహజ సంఖ్యలతో మొదలుపెట్టాం. తరువాత పూర్ణ సంఖ్యలను, అలాగే
``పూర్ణ సంఖ్యల $0$''ను, అంటే $ \equivrep{0,0}{\Intequiv}$ను నిర్మించాం.
కానీ $0 \neq \equivrep{0,0}{\Intequiv}$. నిజానికి మన నిర్మాణాల ప్రకారం
\emph{ఏ} సహజ సంఖ్యా పూర్ణ సంఖ్య కాదు. ఇది అంతర్బుద్ధికి తీవ్రంగా
విరుద్ధంగా అనిపిస్తుంది. నిజానికి \olref[sfr][set][imp]{sec}లో పెద్దగా
వాదించకుండానే $\Nat \subseteq \Rat$ అని చెప్పాం. నిర్మాణాలు సంఖ్యలు
ఖచ్చితంగా \emph{ఏమిటో} చెబుతాయని అనుకుంటే, ఆ వాదన స్పష్టంగానే అసత్యం.

ఒక అడుగు వెనక్కి వేసి చూస్తే, మనం ఎక్కడికి వచ్చాం? సరళ సమితి
సిద్ధాంతంలో పనిచేస్తూ, సహజ సంఖ్యలు మనకు ఇచ్చివున్నాయని తీసుకుని,
పూర్ణ, కరణీయ, వాస్తవ సంఖ్యలను కొన్ని సమితులుగా \emph{పరిగణించగలిగాం}.
ఆ అర్థంలో ఈ వస్తువుల సిద్ధాంతాలను సమితి సిద్ధాంతంలో
\emph{అంతఃస్థాపించగలం}. కానీ ఈ అంతఃస్థాపనకు ఉన్న తాత్త్విక ప్రాముఖ్యత
అంత సూటిగా లేదు.

అయితే, దీనితో చర్చ ముగియదు!{} చెప్పదలచిన విషయం ఇంతే. వాస్తవ సంఖ్యల
అంకగణితీకరణకు \emph{లోతైన} తాత్త్విక ప్రాముఖ్యత ఉందని చూపాలంటే,
మరొక అదనపు \emph{తాత్త్విక} వాదన అవసరం.

%అదనంగా: ఈ అధ్యాయం అంతటా సహజ సంఖ్యలు మనకు సూటిగా
%emph{ఇచ్చివున్నాయని} పరిగణించాం. కానీ వాటితో మనం ఏమి చేయగలం?
%తదుపరి అధ్యాయం చెప్పే విషయం అదే.

\end{document}
